% !TEX program = xelatex
% ============================================================
% SLB 技术文档：6.5.1–6.5.3 共性基础算法（CRC / Gold / QPSK）
% 规范：slb-doc-generator / tech-doc-style-guide.md
% 模块类型：通用功能
% ============================================================
\documentclass[11pt,a4paper]{ctexart}

\usepackage{amsmath,amssymb,amsfonts}
\usepackage{booktabs}
\usepackage{geometry}
\usepackage{hyperref}
\usepackage{enumitem}
\usepackage{xcolor}
\usepackage{longtable}
\usepackage{array}

\geometry{left=2.5cm,right=2.5cm,top=2.5cm,bottom=2.5cm}
\hypersetup{colorlinks=true,linkcolor=blue,urlcolor=blue}

\newcolumntype{L}[1]{>{\raggedright\arraybackslash}p{#1}}

\title{%
  \textbf{星闪 SLB 物理层}\\[0.4em]
  \Large 6.5.1\textasciitilde 6.5.3 共性基础算法技术文档\\[0.3em]
  \large 循环冗余校验 \textbullet\ 伪随机序列 \textbullet\ 伪随机 QPSK 序列
}
\author{依据 T/XS 10001-2025《星闪无线通信系统 接入层\\
同步低时延宽带空口 SLB 技术要求和测试方法》整理\\
（团体标准 20250326 版）}
\date{2026 年 7 月 27 日}

\begin{document}
\maketitle
\tableofcontents
\newpage

%======================================================================
\part{总览}
%======================================================================

\section{文档范围与模块划分}

本文档对应 T/XS 10001-2025 第 6.5.1\textasciitilde 6.5.3 节，覆盖物理层\textbf{共性基础算法}中的循环冗余校验（CRC）、31 阶 Gold 伪随机序列与伪随机 QPSK 参考序列。三节不规定业务语义，只规定``怎么算''比特校验与序列；工程上合并为单一可复用库（\texttt{slb\_phy\_common}），供控制、数据、参考信号与加扰等多条链路调用。

\begin{longtable}{L{2.2cm}L{4.2cm}L{7.5cm}}
\caption{文档范围与模块划分\label{tab:6513-scope}} \\
\toprule
\textbf{子节号} & \textbf{子节主题} & \textbf{核心输出/行为} \\
\midrule
6.5.1 & 循环冗余校验 CRC & 五种 $g(D)$ 编码/校验；系统码 $B{=}A{+}L$；可选 $L{=}24$ Phy-ID 掩码 \\
6.5.2 & 伪随机 Gold 序列 & 双 $m$ 序列 LFSR + 1600 偏移，输出 $c(0)\ldots c(M_{\mathrm{PN}}{-}1)$ \\
6.5.3 & 伪随机 QPSK 序列 & $c_{\mathrm{init}}(n,l,N_{\mathrm{ID}},N_{\mathrm{CP}})$ $\rightarrow$ 161 点 QPSK；$k{=}80$ DC 置零 \\
\bottomrule
\end{longtable}

\section{与相关章节的关系}

\begin{longtable}{L{2.2cm}L{1.8cm}L{9.5cm}}
\caption{与相关章节的关系\label{tab:6513-rel}} \\
\toprule
\textbf{相关节号} & \textbf{关系} & \textbf{说明} \\
\midrule
6.1 & 上游 & 比特序、物理层服务与链路命名约定 \\
6.2.2 & 上游 & 无线帧号 $n$、符号号 $l$、五种 CP-OFDM 类型与 $N_{\mathrm{CP}}$、$a_{80,l}{\equiv}0$ \\
6.2.3 & 下游 & SRS/CSI-RS/DMRS/PAS 等调用 6.5.2+6.5.3 生成参考符号并映射 RE \\
6.2.4 & 上游/下游 & 控制信息比特组装后调 6.5.1；部分 ACK 序列方式调 6.5.3 \\
6.2.5 & 上游/下游 & 数据 TB 信息比特调 6.5.1；第一类数据后续加扰依赖 6.5.2（经 6.5.7） \\
6.2.6 & 下游 & 极化编码输入为 CRC 后比特流 $b_0\ldots b_{B-1}$；码块分段 CRC 复用 6.5.1 \\
6.5.7 & 下游 & 比特加扰：$c_{\mathrm{init}}=n_f\cdot 2^{24}+N_{\mathrm{G-PID}}$，序列由 6.5.2 生成 \\
6.5.9.1 & 下游 & ZC 序列组跳 $c_{\mathrm{init}}=\lfloor n_{\mathrm{ID}}^X/30\rfloor$，复用 6.5.2 \\
6.3 & 平行 & 灵活带宽下同算法，参数随 $\alpha$ 缩放时由调用方调整 $M_{\mathrm{PN}}$/映射 \\
\bottomrule
\end{longtable}

\section{通用功能端到端流水线}

三节构成 PHY``校验—伪随机—参考符号''共性底座。端到端按链路分流，共用同一库、不同参数：

\begin{center}
\small
\textbf{控制面}：6.2.4 组装 $a_0\ldots a_{A-1}$
$\rightarrow$ \textbf{6.5.1} CRC
$\rightarrow$ 6.2.6 极化\\[0.35em]
\textbf{参考信号}：$(n,l,N_{\mathrm{ID}},N_{\mathrm{CP}})$
$\rightarrow$ $c_{\mathrm{init}}$
$\rightarrow$ \textbf{6.5.2} Gold（$M_{\mathrm{PN}}{=}322$）
$\rightarrow$ \textbf{6.5.3} QPSK
$\rightarrow$ 6.2.3 RE 映射（$k{=}80$ 置零）\\[0.35em]
\textbf{数据面}：6.2.5 TB
$\rightarrow$ \textbf{6.5.1} TB/码块 CRC
$\rightarrow$（可选）6.5.7 加扰（\textbf{6.5.2} 提供 $c(i)$）
$\rightarrow$ 编码/调制映射
\end{center}

推荐工程策略：只维护一份 \texttt{slb\_phy\_common}；链路差异体现在 $A$、$g(D)$、$c_{\mathrm{init}}$、$N_{\mathrm{ID}}$ 等入参，不复制算法副本。

\newpage
%======================================================================
\part{6.5.1\textasciitilde 6.5.3 共性基础算法（工程解读）}
%======================================================================

\section{协议章节对应}

本节对应 T/XS 10001-2025 第 6 章第 6.5.1、6.5.2、6.5.3 节，涵盖：
\begin{itemize}[nosep]
  \item \textbf{6.5.1} 循环冗余校验——CRC24A/CRC24B/CRC16/CRC8/CRC6；
  \item \textbf{6.5.2} 伪随机序列——31 阶 Gold（双 $m$ 序列 + 1600 偏移）；
  \item \textbf{6.5.3} 伪随机 QPSK 序列——基于 Gold 的 161 子载波参考符号。
\end{itemize}

\section{功能定义}

\begin{enumerate}[nosep]
  \item \textbf{生成} CRC 校验比特：对信息比特 $a_0\ldots a_{A-1}$ 按选定 $g(D)$ 在 GF(2) 上整除，输出系统码 $b_0\ldots b_{B-1}$（$B{=}A{+}L$）；$L{=}24$ 时可叠加 Phy-ID 掩码；
  \item \textbf{校验} 接收比特流：按相同 $g(D)$/掩码判定通过或失败，仅检测差错，不纠错；
  \item \textbf{生成} Gold 伪随机比特：$c_{\mathrm{init}}$、$M_{\mathrm{PN}}$ $\rightarrow$ $c(0)\ldots c(M_{\mathrm{PN}}{-}1)$；
  \item \textbf{生成} 伪随机 QPSK 参考序列：$n,l,N_{\mathrm{ID}},N_{\mathrm{CP}}$ $\rightarrow$ $r_{n,l}(0)\ldots r_{n,l}(160)$，供参考信号/序列 ACK 映射；
  \item \textbf{映射辅助}：按子载波索引 $k$ 输出 RE 值，$k{=}80$ 强制 $(0,0)$。
\end{enumerate}

\section{模块边界（上下游及职责划分）}

\subsection{上游模块}

\begin{longtable}{L{4.2cm}L{4.5cm}L{5.5cm}}
\caption{上游模块输入\label{tab:6513-up}} \\
\toprule
\textbf{上游} & \textbf{输入} & \textbf{说明} \\
\midrule
控制信息组装（6.2.4） & $a_0\ldots a_{A-1}$ & LSB$\rightarrow$MSB 组装，预留位恒 0 也参与 CRC \\
数据 TB 处理（6.2.5） & TB/码块信息比特 + 分段参数 & $A$ 随 TB/码块变化 \\
帧配置/调度 & CRC 类型、$c_{\mathrm{init}}$ 参数、$M_{\mathrm{PN}}$ & 决定 $g(D)$ 与 Gold 长度 \\
高层 ID 配置 & Phy-ID 掩码（24 bit） & 仅 $L{=}24$ 场景 \\
时频/小区参数 & $n,l,N_{\mathrm{ID}},N_{\mathrm{CP}}$ & 供 6.5.3 计算 $c_{\mathrm{init}}$ \\
\bottomrule
\end{longtable}

\subsection{下游模块}

\begin{longtable}{L{4.2cm}L{4.5cm}L{5.5cm}}
\caption{下游模块输出\label{tab:6513-down}} \\
\toprule
\textbf{下游} & \textbf{本模块输出} & \textbf{说明} \\
\midrule
极化编码（6.2.6） & $b_0\ldots b_{B-1}$ & 控制/码块 CRC 后比特 \\
比特加扰（6.5.7） & Gold $c(i)$ & 与数据比特逐位模 2 加 \\
资源映射（6.2.3） & QPSK $r_{n,l}(k)$ & 161 点；$k{=}80$ 置零后写入网格 \\
序列组跳（6.5.9.1） & Gold $c(i)$ & $c_{\mathrm{init}}=\lfloor n_{\mathrm{ID}}^X/30\rfloor$ \\
HARQ/上层 & CRC 通过标志 & 仅检测；重传决策不在本模块 \\
\bottomrule
\end{longtable}

\subsection{本模块不负责的内容}

\begin{longtable}{L{5.5cm}L{8.5cm}}
\caption{本模块不负责的内容\label{tab:6513-out}} \\
\toprule
\textbf{不负责内容} & \textbf{原因 / 负责节} \\
\midrule
业务字段语义与比特组装顺序细则 & 由 6.2.4/6.2.5 各子节规定；本模块只吃已组装 $a[]$ \\
信道编码（极化/LDPC）与速率匹配 & 6.2.6 \\
资源映射与时频分配策略 & 6.2.3 \\
加扰 $c_{\mathrm{init}}$ 业务公式以外的链路策略 & 6.5.7（本模块只生成 $c(i)$） \\
纠错或 HARQ 重传决策 & CRC 仅检测；上层/MAC \\
\bottomrule
\end{longtable}

\section{在 SLB 通信流程中的角色}

\begin{itemize}[nosep]
  \item \textbf{控制信息链路}：组装 $\rightarrow$ 6.5.1 CRC $\rightarrow$ 6.2.6 极化 $\rightarrow$ 映射；
  \item \textbf{参考信号链路}：调度给出 $(n,l,N_{\mathrm{ID}},N_{\mathrm{CP}})$ $\rightarrow$ 6.5.2/6.5.3 $\rightarrow$ 6.2.3 网格；
  \item \textbf{数据链路}：TB CRC（6.5.1）$\rightarrow$ 可选加扰（6.5.7 调 6.5.2）$\rightarrow$ 编码调制；
  \item \textbf{接收侧}：对控制/数据做 CRC 校验；对参考信号按相同参数本地再生做信道估计。
\end{itemize}

%----------------------------------------------------------------------
\section{关键参数与强制要求}
%----------------------------------------------------------------------

\subsection{6.5.1 循环冗余校验}

\textbf{强制规则}：
\begin{itemize}[nosep]
  \item GF(2) 整除：$T(D)$ 能被所选 $g(D)$ 整除；
  \item 系统码：$b_k=a_k$（$k{=}0\ldots A{-}1$），无掩码时 $b_k=p_{k-A}$（$k{=}A\ldots A{+}L{-}1$）；
  \item 带掩码（仅 $L{=}24$）：$b_k=(p_{k-A}+x_{\mathrm{MASK},k-A})\bmod 2$；
  \item 比特顺序：下标 0 对应 LSB 字段先组装；预留位恒 0 也参与计算。
\end{itemize}

\begin{longtable}{L{2.6cm}L{7.2cm}L{1.2cm}L{3.2cm}}
\caption{CRC 生成多项式与典型用途\label{tab:6513-crc-poly}} \\
\toprule
\textbf{枚举} & \textbf{生成多项式 $g(D)$} & \textbf{$L$} & \textbf{120\,kHz 典型用途} \\
\midrule
\texttt{SLB\_CRC24A} &
$D^{24}{+}D^{23}{+}D^{18}{+}D^{17}{+}D^{14}{+}D^{11}{+}D^{10}{+}D^{7}{+}D^{6}{+}D^{5}{+}D^{4}{+}D^{3}{+}D{+}1$ &
24 & 第二类数据 TB \\
\texttt{SLB\_CRC24B} &
$D^{24}{+}D^{23}{+}D^{21}{+}D^{20}{+}D^{17}{+}D^{15}{+}D^{13}{+}D^{12}{+}D^{8}{+}D^{4}{+}D^{2}{+}D{+}1$ &
24 & SAB、HARQ-ACK、GCI、极化分段 \\
\texttt{SLB\_CRC16} & $D^{16}{+}D^{12}{+}D^{5}{+}1$ & 16 & CQI、调度请求 SR \\
\texttt{SLB\_CRC8} & $D^{8}{+}D^{6}{+}D^{5}{+}D^{3}{+}1$ & 8 & 第一类数据可选 \\
\texttt{SLB\_CRC6} & $D^{6}{+}D^{5}{+}1$ & 6 & 特殊短控制信息 \\
\bottomrule
\end{longtable}

\subsection{6.5.2 伪随机 Gold 序列}

\begin{align}
c(i) &= \bigl(x_1(i{+}1600) + x_2(i{+}1600)\bigr) \bmod 2 \\
x_1(i{+}31) &= \bigl(x_1(i{+}3) + x_1(i)\bigr) \bmod 2 \\
x_2(i{+}31) &= \bigl(x_2(i{+}3) + x_2(i{+}2) + x_2(i{+}1) + x_2(i)\bigr) \bmod 2
\end{align}

初始化：$x_1(0){=}1$，$x_1(i){=}0$（$i{=}1\ldots 30$）；$c_{\mathrm{init}}=\sum_{i=0}^{30} x_2(i)\cdot 2^i$（LSB 对应 $i{=}0$）。

\textbf{1600 偏移}：输出取 $x_1(i{+}1600)$、$x_2(i{+}1600)$，跳过 $m$ 序列瞬态（与 3GPP 同类约定一致）；禁止省略偏移。

\textbf{$x_1$ 预计算优化}：$x_1$ 初态固定、与 $c_{\mathrm{init}}$ 无关，可将 $x_1(0)\ldots x_1(1600{+}M_{\mathrm{PN}}{-}1)$ 固化为只读表；运行时仅递推 $x_2$ 再异或，输出与逐步递推等价。

\begin{longtable}{L{4.5cm}L{2.2cm}L{7.5cm}}
\caption{典型链路 $c_{\mathrm{init}}$ 公式（调用方计算后传入）\label{tab:6513-cinit}} \\
\toprule
\textbf{场景} & \textbf{标准节} & \textbf{$c_{\mathrm{init}}$ 公式} \\
\midrule
伪随机 QPSK 参考信号 & 6.5.3 &
$2^{10}(15(n{+}1){+}l{+}1)(4N_{\mathrm{ID}}{+}1){+}4N_{\mathrm{ID}}{+}N_{\mathrm{CP}}$ \\
比特加扰 & 6.5.7 & $n_f\cdot 2^{24} + N_{\mathrm{G-PID}}$ \\
ZC 序列组跳 & 6.5.9.1 & $\lfloor n_{\mathrm{ID}}^X / 30\rfloor$ \\
\bottomrule
\end{longtable}

\subsection{6.5.3 伪随机 QPSK 序列}

\begin{equation}
r_{n,l}(m) = \frac{1}{\sqrt{2}}\bigl(2c(2m)-1\bigr) + j\,\frac{1}{\sqrt{2}}\bigl(2c(2m{+}1)-1\bigr), \quad m=0,\ldots,160
\end{equation}

所需 Gold 长度 $M_{\mathrm{PN}}=322$；$c_{\mathrm{init}}$ 见上表 6.5.3 行。资源映射：$k{\neq}80$ 时 $a_{k,l}=r_{n,l}(k)$，$k{=}80$ 置零。

\begin{longtable}{L{2.5cm}L{4.5cm}L{7cm}}
\caption{6.5.3 $c_{\mathrm{init}}$ 参数含义\label{tab:6513-qpsk-param}} \\
\toprule
\textbf{参数} & \textbf{含义} & \textbf{取值范围} \\
\midrule
$n$ & 超帧内无线帧号 & $0,1,\ldots,7$ \\
$l$ & CP-OFDM 符号编号 & 随符号类型变化（见表~\ref{tab:6513-ncp}） \\
$N_{\mathrm{ID}}$ & G 节点标识低 8 位 & $0,1,\ldots,255$ \\
$N_{\mathrm{CP}}$ & 循环前缀类型编码 & $0,1,2,3$ \\
\bottomrule
\end{longtable}

\begin{longtable}{L{4.8cm}L{3.8cm}L{1.8cm}L{3.6cm}}
\caption{CP-OFDM 符号类型与 $l$、$N_{\mathrm{CP}}$ 对应\label{tab:6513-ncp}} \\
\toprule
\textbf{CP-OFDM 符号类型} & \textbf{符号编号 $l$} & \textbf{$N_{\mathrm{CP}}$} & \textbf{说明} \\
\midrule
类型一 A 或类型一 B & $0\ldots 13$ & 3 & 每无线帧 14 符号 \\
类型二 & $0\ldots 12$ & 2 & 每无线帧 13 符号 \\
类型三 & $0\ldots 11$ & 1 & 每无线帧 12 符号 \\
类型四 & $0\ldots 9$ & 0 & 每无线帧 10 符号 \\
\bottomrule
\end{longtable}

\textbf{标准算例}（$n{=}0,l{=}0,N_{\mathrm{ID}}{=}5,N_{\mathrm{CP}}{=}3$）：
$c_{\mathrm{init}}=344087$，$c(0\ldots 4)=0,1,1,1,0$。联调应以该向量作 Gold/QPSK 自检金标。

\begin{longtable}{cccL{6cm}}
\caption{QPSK 比特对到星座点映射\label{tab:6513-qpsk-map}} \\
\toprule
\textbf{$c(2m)$} & \textbf{$c(2m{+}1)$} & \textbf{$r_{n,l}(m)$} & \textbf{说明} \\
\midrule
0 & 0 & $(-1-j)/\sqrt{2}$ & I/Q 同为负 \\
0 & 1 & $(-1+j)/\sqrt{2}$ & \\
1 & 0 & $(+1-j)/\sqrt{2}$ & \\
1 & 1 & $(+1+j)/\sqrt{2}$ & \\
\bottomrule
\end{longtable}

%----------------------------------------------------------------------
\section{本节核心详解}
%----------------------------------------------------------------------

\subsection{6.5.1 CRC 编码与校验}

\textbf{功能}：对已组装信息比特附加/核验 CRC。\\
\textbf{输入/输出}：见接口表~\ref{tab:6513-if-crc}；编码输出长度 $B{=}A{+}L$。\\
\textbf{上下文}：控制面在极化前；数据面在 TB/码块级；接收面在译码后判通过。

掩码仅用于标准明示的 $L{=}24$ 场景（如部分 HARQ-ACK/GCI/预配置调度激活去激活）；$L{\neq}24$ 时 \texttt{mask} 必须为空。

\subsection{6.5.2 Gold 序列生成}

\textbf{功能}：按 $c_{\mathrm{init}}$ 生成任意长度伪随机比特。\\
\textbf{输入/输出}：$c_{\mathrm{init}}$、$M_{\mathrm{PN}}$ $\rightarrow$ $c[]$。\\
\textbf{上下文}：6.5.3（固定 322）、6.5.7 加扰、6.5.9.1 组跳；调用方负责业务侧 $c_{\mathrm{init}}$ 公式。

\subsection{6.5.3 QPSK 参考序列与 RE 映射}

\textbf{功能}：由时频/小区参数生成 161 点单位能量 QPSK，并按 $k$ 写入 RE。\\
\textbf{输入/输出}：$(n,l,N_{\mathrm{ID}},N_{\mathrm{CP}})$ $\rightarrow$ $r[161]$；$k$ $\rightarrow$ $a_{k,l}$。\\
\textbf{上下文}：SRS/CSI-RS/DMRS/PAS 及 T 链路序列 ACK；每符号 $l$ 可独立再生。

%----------------------------------------------------------------------
\section{数据类型、长度与维度}
%----------------------------------------------------------------------

\begin{longtable}{L{3.2cm}L{2.8cm}L{2.2cm}L{6cm}}
\caption{CRC 编码接口维度（\texttt{slb\_crc\_encode}）\label{tab:6513-dim-crc}} \\
\toprule
\textbf{变量} & \textbf{方向} & \textbf{长度} & \textbf{说明} \\
\midrule
\texttt{a} & 输入 & $A$ bit & 信息比特，元素 0/1 \\
\texttt{A} & 输入 & 标量 & 信息比特个数 \\
\texttt{poly} & 输入 & 枚举 & 五种多项式之一 \\
\texttt{mask} & 输入 & 24 bit 或空 & 仅 $L{=}24$；空表示不加掩码 \\
\texttt{b} & 输出 & $A{+}L$ bit & 编码比特 \\
\texttt{b\_cap} & 输入 & 标量 & 须 $\geq A{+}L$ \\
\bottomrule
\end{longtable}

\begin{longtable}{L{3.2cm}L{2.8cm}L{2.2cm}L{6cm}}
\caption{Gold 序列接口维度（\texttt{slb\_pn\_gold\_generate}）\label{tab:6513-dim-gold}} \\
\toprule
\textbf{变量} & \textbf{方向} & \textbf{长度} & \textbf{说明} \\
\midrule
\texttt{c\_init} & 输入 & $\le 31$ bit 有效 & Gold 初始化整数 \\
\texttt{M\_PN} & 输入 & 标量 & 所需序列长度 \\
\texttt{c} & 输出 & $M_{\mathrm{PN}}$ & $c(0)\ldots c(M_{\mathrm{PN}}{-}1)$ \\
\texttt{c\_cap} & 输入 & 标量 & 须 $\geq M_{\mathrm{PN}}$ \\
\bottomrule
\end{longtable}

\begin{longtable}{L{3.2cm}L{2.8cm}L{2.2cm}L{6cm}}
\caption{QPSK 接口维度\label{tab:6513-dim-qpsk}} \\
\toprule
\textbf{变量} & \textbf{方向} & \textbf{长度} & \textbf{说明} \\
\midrule
\texttt{param.n} & 输入 & 0\ldots 7 & 无线帧号 \\
\texttt{param.l} & 输入 & 见类型表 & 符号号 \\
\texttt{param.N\_ID} & 输入 & 0\ldots 255 & G-ID 低 8 位 \\
\texttt{param.N\_CP} & 输入 & 0\ldots 3 & CP 类型编码 \\
\texttt{r} & 输出 & 161 复数 & $r_{n,l}(0)\ldots r_{n,l}(160)$ \\
\texttt{r\_cap} & 输入 & 标量 & 须 $\geq 161$ \\
\bottomrule
\end{longtable}

%----------------------------------------------------------------------
\section{时序关系与触发条件}
%----------------------------------------------------------------------

\begin{longtable}{L{3.5cm}L{5.5cm}L{5.5cm}}
\caption{调用触发条件\label{tab:6513-trig}} \\
\toprule
\textbf{能力} & \textbf{进入条件} & \textbf{失败/回退} \\
\midrule
CRC 编码 & 信息比特已按 6.2.4/6.2.5 组装完成；$A$、$g(D)$ 已知 & 空指针/容量不足/非法 poly $\rightarrow$ 返回错误，不写半结果 \\
CRC 校验 & 译码输出比特流与发送侧 poly/掩码一致 & 不通过 $\rightarrow$ 上报失败，本模块不触发重传 \\
Gold 生成 & $c_{\mathrm{init}}$、$M_{\mathrm{PN}}$ 已由上层算好 & 容量不足 $\rightarrow$ 错误；禁止无 1600 偏移的``快路径'' \\
QPSK 生成 & 帧号/符号号/ID/CP 类型有效 & $l$ 超出当前类型范围 $\rightarrow$ 参数错误 \\
RE 映射 & $r[]$ 已生成；$k\in[0,160]$ & $k{=}80$ 输出 $(0,0)$，不视为错误 \\
\bottomrule
\end{longtable}

\textbf{工作模式}：纯函数调用；无持续流式内部状态。\\
\textbf{实时性}：CRC 为 $O(A)$；Gold 对 $x_2$ 需 $M_{\mathrm{PN}}{+}1600$ 次递推，建议帧起始前预计算或查表 $x_1$。

%----------------------------------------------------------------------
\section{接口表格（明确的输入/输出）}
%----------------------------------------------------------------------

\begin{longtable}{L{2.2cm}L{3.0cm}L{1.2cm}L{1.5cm}L{2.2cm}L{4.2cm}}
\caption{\texttt{slb\_crc\_encode} 接口\label{tab:6513-if-crc}} \\
\toprule
\textbf{参数} & \textbf{类型} & \textbf{长度} & \textbf{必需} & \textbf{来源} & \textbf{说明} \\
\midrule
\texttt{a} & \texttt{const uint8\_t *} & $A$ & 是 & 信息组装 & 比特 0/1 \\
\texttt{A} & \texttt{uint32\_t} & 1 & 是 & 上游 & 信息比特数 \\
\texttt{poly} & \texttt{slb\_crc\_poly\_t} & 1 & 是 & 协议配置 & CRC 类型 \\
\texttt{mask} & \texttt{const uint8\_t *} & 24 & 条件 & 高层 ID & $L{=}24$ 掩码 \\
\texttt{b} & \texttt{uint8\_t *} & $A{+}L$ & 是 & 调用方 & 输出缓冲 \\
\texttt{b\_cap} & \texttt{uint32\_t} & 1 & 是 & 调用方 & $\geq A{+}L$ \\
\bottomrule
\end{longtable}

\begin{longtable}{L{3.5cm}L{10cm}}
\caption{其余 API 摘要\label{tab:6513-if-rest}} \\
\toprule
\textbf{API} & \textbf{输入 / 输出要点} \\
\midrule
\texttt{slb\_crc\_check} & 入：\texttt{b}（$B$ bit）、\texttt{poly}、\texttt{mask}；出：\texttt{pass}$\in\{0,1\}$ \\
\texttt{slb\_pn\_gold\_generate} & 入：\texttt{c\_init}、\texttt{M\_PN}、\texttt{c\_cap}；出：\texttt{c[0..M\_PN-1]} \\
\texttt{slb\_qpsk\_calc\_cinit} & 入：\texttt{n,l,N\_ID,N\_CP}；出：$c_{\mathrm{init}}$ \\
\texttt{slb\_qpsk\_ref\_generate} & 内调 Gold（322）+ QPSK 映射；出：161 复数 \\
\texttt{slb\_qpsk\_map\_to\_re} & 入：$r[]$、$k$；出：$a_{k,l}$（$k{=}80$ 为 0） \\
\bottomrule
\end{longtable}

\begin{longtable}{L{4cm}L{2cm}L{8cm}}
\caption{错误码\label{tab:6513-err}} \\
\toprule
\textbf{宏} & \textbf{值} & \textbf{含义} \\
\midrule
\texttt{SLB\_PHY\_OK} & 0 & 成功 \\
\texttt{SLB\_PHY\_ERR\_PARAM} & $-1$ & 空指针或非法参数 \\
\texttt{SLB\_PHY\_ERR\_BUF} & $-2$ & 输出缓冲区容量不足 \\
\bottomrule
\end{longtable}

%----------------------------------------------------------------------
\section{处理步骤}
%----------------------------------------------------------------------

\subsection{端到端链路 A：控制信息 CRC $\rightarrow$ 极化}

\textbf{Step 1（组装完成 $\rightarrow$ 选多项式）}
\begin{itemize}[nosep]
  \item \textbf{进入}：$a_0\ldots a_{A-1}$ 已按 6.2.4 组装（含预留 0）；
  \item \textbf{动作}：按信息块类型查表选 $g(D)$（如 SAB$\rightarrow$CRC24B，CQI$\rightarrow$CRC16）；若标准要求掩码则准备 24 bit Phy-ID；
  \item \textbf{失败}：$A$ 与字段定义不符或 poly 非法 $\rightarrow$ \texttt{SLB\_PHY\_ERR\_PARAM}。
\end{itemize}

\textbf{Step 2（编码）}
\begin{itemize}[nosep]
  \item \textbf{进入}：参数合法且 \texttt{b\_cap}$\geq A{+}L$；
  \item \textbf{动作}：CRC 寄存器清零；逐 bit 移入 $a_i$ 按抽头反馈；取 $p_0\ldots p_{L-1}$；有掩码则 $b_{A+i}=(p_i+\mathrm{mask}_i)\bmod 2$，否则 $b_{A+i}=p_i$；前 $A$ 位拷贝 $a$；
  \item \textbf{失败}：容量不足 $\rightarrow$ \texttt{SLB\_PHY\_ERR\_BUF}，不部分写回。
\end{itemize}

\textbf{Step 3（交给极化）}
\begin{itemize}[nosep]
  \item \textbf{进入}：得到 $b_0\ldots b_{B-1}$；
  \item \textbf{动作}：作为 6.2.6 输入；本模块结束；
  \item \textbf{失败}：若下游发现长度不等于 $A{+}L$，回查组装/$L$，不在 CRC 内``补丁改长''。
\end{itemize}

\subsection{端到端链路 B：参考信号 Gold $\rightarrow$ QPSK $\rightarrow$ RE}

\textbf{Step 1（参数 $\rightarrow$ $c_{\mathrm{init}}$）}
\begin{itemize}[nosep]
  \item \textbf{进入}：$n,l,N_{\mathrm{ID}},N_{\mathrm{CP}}$ 有效且 $l$ 落在当前 CP 类型允许范围；
  \item \textbf{动作}：$c_{\mathrm{init}}=2^{10}(15(n{+}1)+l{+}1)(4N_{\mathrm{ID}}{+}1)+4N_{\mathrm{ID}}+N_{\mathrm{CP}}$；
  \item \textbf{失败}：越界参数 $\rightarrow$ 返回错误，不生成序列。
\end{itemize}

\textbf{Step 2（Gold，$M_{\mathrm{PN}}{=}322$）}
\begin{itemize}[nosep]
  \item \textbf{进入}：$c_{\mathrm{init}}$ 已算；
  \item \textbf{动作}：初始化 $x_1$、$x_2$；递推至 $1600{+}M_{\mathrm{PN}}$；输出 $c(i)=(x_1(i{+}1600)+x_2(i{+}1600))\bmod 2$（可用 $x_1$ 查表）；
  \item \textbf{失败}：省略 1600 偏移的实现判为不合格。
\end{itemize}

\textbf{Step 3（QPSK 映射）}
\begin{itemize}[nosep]
  \item \textbf{进入}：$c[0\ldots 321]$ 就绪；
  \item \textbf{动作}：对 $m{=}0\ldots 160$ 按表~\ref{tab:6513-qpsk-map} 生成 $r_{n,l}(m)$；
  \item \textbf{失败}：\texttt{r\_cap}${<}161$ $\rightarrow$ \texttt{SLB\_PHY\_ERR\_BUF}。
\end{itemize}

\textbf{Step 4（RE 写入）}
\begin{itemize}[nosep]
  \item \textbf{进入}：调用方遍历 $k{=}0\ldots 160$；
  \item \textbf{动作}：$k{\neq}80$ 写 $r_{n,l}(k)$，$k{=}80$ 写 $(0,0)$；
  \item \textbf{失败}：把 DC 当有效参考符号使用 $\rightarrow$ 信道估计联调失败。
\end{itemize}

\subsection{端到端链路 C：数据 TB CRC +（可选）加扰}

\textbf{Step 1}：按 6.2.5 选定 CRC24A/16/8 等对 TB/码块编码（同链路 A 的 Step 2）。\\
\textbf{Step 2}：若链路需加扰，按 6.5.7 算 $c_{\mathrm{init}}=n_f\cdot 2^{24}+N_{\mathrm{G-PID}}$，调 6.5.2 生成 $c(i)$，再 $\tilde{t}_i=(t_i+c(i))\bmod 2$。\\
\textbf{Step 3}：进入编码/调制映射；本模块不参与 MCS/RE 分配。

\subsection{接收侧 CRC 校验}

\begin{itemize}[nosep]
  \item \textbf{进入}：译码得到候选比特流与已知 poly/掩码；
  \item \textbf{动作}：若有掩码先还原校验位再整除检验，或等价地按发送规则重算比较；输出 \texttt{pass}；
  \item \textbf{失败}：\texttt{pass}{=}0 时仅上报，由上层决定丢弃/重传。
\end{itemize}

%----------------------------------------------------------------------
\section{协议中允许本模块改变的参数}
%----------------------------------------------------------------------

\begin{longtable}{L{4cm}L{4cm}L{6cm}}
\caption{允许配置的参数及其影响\label{tab:6513-tunable}} \\
\toprule
\textbf{参数} & \textbf{来源} & \textbf{影响} \\
\midrule
\texttt{poly} & 帧类型/信令 & 决定 $L$ 与 $g(D)$ \\
\texttt{A} & 上游输入 & 改变 $B{=}A{+}L$ \\
\texttt{mask} & 高层 ID & 仅 $L{=}24$，改变校验位 \\
\texttt{c\_init}, \texttt{M\_PN} & 调度/ID 配置 & 改变 Gold/QPSK 序列 \\
\texttt{n,l,N\_ID,N\_CP} & 帧/CP 配置 & 改变 QPSK 的 $c_{\mathrm{init}}$ \\
\bottomrule
\end{longtable}

%----------------------------------------------------------------------
\section{链路调用对照}
%----------------------------------------------------------------------

\begin{longtable}{L{4.2cm}L{2.8cm}ccc}
\caption{链路与 6.5.1/6.5.2/6.5.3 调用对照\label{tab:6513-link}} \\
\toprule
\textbf{链路/信息块} & \textbf{标准位置} & \textbf{6.5.1} & \textbf{6.5.2} & \textbf{6.5.3} \\
\midrule
通信域 SAB & 6.2.4.1.1 & CRC24B & --- & --- \\
直通同步信息 & 6.2.4.1.1 & CRC24B & --- & --- \\
HARQ-ACK / GCI & 6.2.4.6/10 & CRC24B$\pm$掩码 & --- & --- \\
CQI / SR & 6.2.4.10.3/4 & CRC16 & --- & --- \\
第二类数据 TB & 6.2.5.3 & CRC24A & --- & --- \\
第一类数据 & 6.2.5.2.2 & CRC24A/16/8 & 经 6.5.7 & --- \\
SRS / CSI-RS / DMRS & 6.2.3 & --- & $c(i)$ & $r_{n,l}(k)$ \\
T 链路 ACK（序列） & 6.2.4.10.6 & --- & $c(i)$ & $\pm r_{n,l}(k)$ \\
极化码块分段 CRC & 6.2.6.1.1 & CRC24B & --- & --- \\
\bottomrule
\end{longtable}

%----------------------------------------------------------------------
\section{约束与实现要点}
%----------------------------------------------------------------------

\begin{enumerate}[nosep]
  \item \textbf{【多项式合法】}：仅允许五种标准 $g(D)$；非法 \texttt{poly} 返回 \texttt{SLB\_PHY\_ERR\_PARAM}；
  \item \textbf{【缓冲容量】}：\texttt{b\_cap}$\geq A{+}L$，\texttt{c\_cap}$\geq M_{\mathrm{PN}}$，\texttt{r\_cap}$\geq 161$，不足返回 \texttt{SLB\_PHY\_ERR\_BUF}；
  \item \textbf{【1600 偏移】}：Gold 必须使用 $x_1(i{+}1600)$/$x_2(i{+}1600)$，禁止从 $i{=}0$ 直接输出；
  \item \textbf{【$x_1$ 预计算】}：$x_1$ 初态固定，建议离线查表；运行时仅递推 $x_2$，结果须与标准逐步递推一致；
  \item \textbf{【DC 置零】}：$k{=}80$ 必须为 $(0,0)$，由映射接口保证，禁止写入有效能量；
  \item \textbf{【比特顺序】}：$a[]/b[]/c[]$ 下标 0 对齐标准公式下标 0（LSB 先组装）；
  \item \textbf{【掩码范围】}：仅 $L{=}24$ 允许掩码；$L{\neq}24$ 传入非空掩码视为参数错误；
  \item \textbf{【可重入】}：无静态全局可变状态，输出缓冲由调用方提供，可多线程并行；
  \item \textbf{【金标自检】}：实现须通过 $n{=}0,l{=}0,N_{\mathrm{ID}}{=}5,N_{\mathrm{CP}}{=}3$ $\Rightarrow$ $c_{\mathrm{init}}{=}344087$，$c(0\ldots 4){=}0,1,1,1,0$；
  \item \textbf{【无状态机】}：CRC/Gold/QPSK 均按``进入条件 $\rightarrow$ 动作 $\rightarrow$ 失败返回''推进，不维护 IDLE/BUSY/DONE 等 FSM 跳转表；失败立即报错或交还调用方，避免状态锁死。
\end{enumerate}

%----------------------------------------------------------------------
\section{与标准其他章节的衔接}
%----------------------------------------------------------------------

\begin{itemize}[nosep]
  \item \textbf{6.1}：比特序与物理层服务约定，决定数组下标语义；
  \item \textbf{6.2.2}：提供 $n$、$l$、CP 类型及 $a_{80,l}{\equiv}0$ 网格约束；
  \item \textbf{6.2.3}：参考信号映射消费 6.5.3 输出；
  \item \textbf{6.2.4/6.2.5}：组装信息比特并规定各信息块选用的 CRC 类型/掩码；
  \item \textbf{6.2.6}：以 CRC 后比特为极化/码块处理输入；分段 CRC 复用 6.5.1；
  \item \textbf{6.5.7}：加扰 $c_{\mathrm{init}}$ 公式在该节，序列生成复用 6.5.2；
  \item \textbf{6.5.9.1}：序列组跳 $c_{\mathrm{init}}$ 复用 6.5.2；
  \item \textbf{6.3}：灵活带宽下算法不变，长度/映射由调用方按 $\alpha$ 调整。
\end{itemize}

\vfill
\begin{center}
\small\textcolor{gray}{精确参数与字段定义以 T/XS 10001-2025 正式标准为准；\\
本文档提供 6.5.1\textasciitilde 6.5.3 共性基础算法工程解读，供 SLB 实现与联调参考。}
\end{center}

\end{document}
